\section{Данные и протокол эксперимента}
Используется англоязычный корпус Databricks Dolly 15k: инструкции и ответы написаны сотрудниками Databricks, а часть контекстов содержит фрагменты Википедии. Корпус распространяется по лицензии CC BY-SA 3.0; авторство принадлежит Databricks и, для соответствующих фрагментов, авторам Википедии. Источник: \url{https://huggingface.co/datasets/databricks/databricks-dolly-15k}.

Каждый пример хранится в JSONL и содержит стабильный идентификатор, категорию, инструкцию, контекст, ответ и метку происхождения. Контекст при отсутствии представлен пустой строкой. Исходные категории сохранены без переименования. Для воспроизводимого разделения используется стратификация по категориям с seed 42: 12000 примеров в train\_pool, 1000 в validation и 2000 в test. Из train\_pool стратифицированно выбираются 2000 реальных обучающих примеров с seed 43. Остаток исходного корпуса не используется. Хеши файлов и исходного набора фиксируются в манифесте. Validation и test состоят только из реальных данных и после утверждения не изменяются.

Seed-задачи для генерации должны совпадать с записями train\_pool. Независимая проверка вычисляет ROUGE-L F1 между каждой синтетической инструкцией и инструкциями в validation и test; совпадения с оценкой не ниже 0,6 не допускаются. Проверка проводится на финальных данных, а её результат вместе с хешами входных файлов сохраняется отдельно. Факт прохождения проверки следует указать только после получения и аудита финального набора.

\textbf{Ограничения.} Категории и тематика отражают особенности исходного корпуса; возможны ошибки и смещения его авторов и Википедии. Разделение по идентификаторам само по себе не исключает близких по смыслу примеров.

В исходных данных встречаются повторяющиеся формулировки: после перевода в нижний регистр и удаления внешних пробелов 12 различных пар «инструкция и контекст» из train\_real совпадают с парами в test (14 тестовых записей), ещё 5 пар совпадают с validation. Эти совпадения могут завышать оценку обобщения; необходимо учитывать их при анализе и согласовать дальнейшие действия с командой, сохраняя зафиксированные сплиты до такого решения.
